\chapter{Ferramentas e métodos}

\section{Ferramentas}

As ferramentas utilizadas neste trabalho foram selecionadas para cobrir análise de qualidade, implementação de testes e versionamento do código no contexto do \textit{AgroMart Web}.

\begin{itemize}
	\item \textbf{SonarQube}: utilizado para análise estática e acompanhamento de métricas de qualidade, com foco em cobertura de testes, \textit{code smells}, duplicidade e indicadores de manutenibilidade~\cite{sonarqube};
	\item \textbf{React Testing Library}: utilizada para implementação de testes unitários com foco no comportamento dos componentes React;
	\item \textbf{Jest}: utilizado como executor de testes, asserções e coleta de cobertura no projeto, integrado ao \textit{react-scripts test};
	\item \textbf{TypeScript}: utilizado no desenvolvimento do \textit{AgroMart Web}, impactando diretamente a escrita e validação dos testes;
	\item \textbf{Git e GitHub}: utilizados para controle de versão, organização das contribuições e revisão das alterações por meio de \textit{pull requests};
	\item \textbf{GitHub Actions}: utilizado para automatizar etapas de integração contínua, como execução de testes e validações a cada \textit{push} ou \textit{pull request}.
\end{itemize}

Além dessas ferramentas, também foram utilizados os scripts nativos do projeto para execução local dos testes e validação incremental das alterações.

\section{Metodologia}

A metodologia adotada neste trabalho é de natureza \textbf{exploratória}, com foco em diagnóstico e melhoria incremental da cobertura de testes no AgroMart Web.
Assim, o estudo não se caracteriza como experimento controlado, pois não há manipulação de grupos, variáveis independentes e dependentes ou comparação estatística entre tratamentos.

A condução metodológica foi organizada em ciclos de exploração e intervenção técnica:

\begin{enumerate}
	\item \textbf{Diagnóstico inicial}: execução de análise global no SonarQube para estabelecer a linha de base das métricas e identificar os problemas prioritários;
	\item \textbf{Mapeamento de unidades fundamentais}: análise do código para identificar unidades de comportamento essenciais ao funcionamento da aplicação;
	\item \textbf{Incremento de testes}: escrita e execução de testes unitários nas unidades priorizadas, com validação local e em \textit{workflow} de \textit{CI};
	\item \textbf{Registro dos resultados}: comparação das métricas antes e depois de cada incremento, especialmente cobertura de testes e manutenibilidade;
	\item \textbf{Consolidação}: sistematização das decisões e práticas em documentação de desenvolvimento para apoiar contribuições futuras.
\end{enumerate}

Dessa forma, os resultados apresentados representam evidências de evolução técnica obtidas durante a exploração do problema e a aplicação de incrementos de teste, e não inferências causais de um desenho experimental.

\section{Cronograma}

A Tabela a seguir apresenta o cronograma planejado para o TCC 2, considerando o período letivo de 10/08 a 14/12/2026.

\begin{table}[h]
\centering
\caption{Cronograma de execução do TCC 2}
\begin{tabular}{|c|p{10cm}|}
\hline
\textbf{Período} & \textbf{Atividade} \\ \hline
10/08 a 23/08 & Planejamento do TCC 2, revisão do escopo e configuração inicial do ambiente de análise e testes. \\ \hline
24/08 a 13/09 & Revisão dos resultados já obtidos no SonarQube, atualização da linha de base e priorização das unidades fundamentais para incremento de testes. \\ \hline
14/09 a 11/10 & Implementação do primeiro ciclo de testes unitários, execução local e em CI, com registro dos resultados parciais de cobertura. \\ \hline
12/10 a 01/11 & Implementação do segundo ciclo de testes unitários, refinamento de casos de teste e reavaliação das métricas de qualidade. \\ \hline
02/11 a 15/11 & Consolidação da documentação de desenvolvimento com diretrizes de testes e organização dos artefatos de evidência técnica. \\ \hline
16/11 a 30/11 & Análise comparativa dos resultados (antes/depois), discussão dos ganhos obtidos e redação dos capítulos finais. \\ \hline
01/12 a 10/12 & Revisão final do texto completo, ajustes de formatação, referências e validação de consistência geral do documento. \\ \hline
11/12 a 14/12 & Fechamento do trabalho, preparação para entrega/defesa e submissão da versão final. \\ \hline
\end{tabular}
\end{table}


